% Part: sets-functions-relations
% Chapter: relations
% Section: operations

\documentclass[../../../include/open-logic-section]{subfiles}

\begin{document}

\olfileid[tr]{sfr}{rel}{ops}
\olsection{Bağıntılar Üzerindeki İşlemler}

Bağıntıları değiştirmek veya birleştirmek çoğu zaman yararlıdır.
\olref[sfr][rel][ord]{prop:stricttopartial} içinde bağıntıların
\emph{birleşimini} ele aldık; bu yalnızca, sıralı ikililerden oluşan kümeler
olarak görülen iki bağıntının birleşimidir. Benzer biçimde
\olref[sfr][rel][ord]{prop:partialtostrict} içinde bağıntılar arasındaki küme farkını
ele aldık. Bağıntılar üzerinde yapabileceğimiz başka işlemler de şunlardır.

\begin{defn}\ollabel{relationoperations}
$R$ ve $S$ bağıntılar, $A$ da herhangi bir küme olsun.

$R$'nin \emph{tersi},
$R^{-1}=\Setabs{\tuple{y,x}}{\tuple{x,y}\in R}$ bağıntısıdır.

$R$ ile $S$'nin \emph{bileşkesi},
$(R \mid S)=\{\tuple{x,z}:\exists y(Rxy\land Syz)\}$ bağıntısıdır.

$R$'nin $A$'ya \emph{kısıtlaması},
$\funrestrictionto{R}{A}=R\cap A^2$ bağıntısıdır.

$R$'nin $A$'ya \emph{uygulanması},
$\funimage{R}{A}=\{y:(\exists x \in A)Rxy\}$ kümesidir.
\end{defn}

\begin{ex}
$S \subseteq \Int^2$, $\Int$ üzerindeki ardıl bağıntısı olsun; yani
$S=\Setabs{\tuple{x,y}\in\Int^2}{x+1=y}$ ve dolayısıyla $Sxy$ ancak ve ancak
$x+1=y$ olduğunda geçerli olsun.

$S^{-1}$, $\Int$ üzerindeki öncül bağıntısıdır; yani
$\Setabs{\tuple{x,y}\in\Int^2}{x-1=y}$ bağıntısıdır.

$S\mid S$, $\Setabs{\tuple{x,y}\in\Int^2}{x+2=y}$ bağıntısıdır.

$\funrestrictionto{S}{\Nat}$, $\Nat$ üzerindeki ardıl bağıntısıdır.

$\funimage{S}{\{1,2,3\}}$, $\{2,3,4\}$ kümesidir.
\end{ex}

\begin{defn}[Geçişli kapanış]
$R \subseteq A^2$ bir ikili bağıntı olsun.

$R$'nin \emph{geçişli kapanışı},
$R^+=\bigcup_{0<n\in\Nat}R^n$ bağıntısıdır; burada $R^1=R$ ve
$R^{n+1}=R^n\mid R$ ifadelerini özyinelemeli olarak tanımlarız.

$R$'nin \emph{yansımalı geçişli kapanışı}, $R^*=R^+\cup\Id{A}$ bağıntısıdır.
\end{defn}

\begin{ex}
$S \subseteq \Int^2$ ardıl bağıntısını ele alalım. $S^2xy$ ancak ve ancak
$x+2=y$, $S^3xy$ ancak ve ancak $x+3=y$ vb. olur. Dolayısıyla $S^+xy$ ancak
ve ancak $x+n=y$ olacak biçimde bir $n\geq 1$ bulunduğunda geçerlidir. Başka bir deyişle
$S^+xy$ ancak ve ancak $x<y$ olduğunda; $S^*xy$ ise ancak ve ancak $x\le y$
olduğunda geçerlidir.
\end{ex}

\begin{prob}
$R$'nin geçişli kapanışının gerçekten geçişli olduğunu gösterin.
\end{prob}

\end{document}
